\chapter*{译\ 者\ 序}
\phantomsection
\addcontentsline{toc}{chapter}{译者序}
\markboth{译者序}{译者序}

这本书到我们手上的时候，是一份打印稿和一只硬盘。

打印稿是日文，二百多页，页边写满了作者自己的批注，有些批注比正文长。硬盘里是同一部稿子的电子版，另加三百来张翻拍的照片，编号混乱，文件名里混着俄文、日文和数字。

原稿是分章交来的。孙笑川译一章，李文亚核一章里的俄文引文，两边的进度始终差着一章。所以头三个月做的是最费事的活：把照片按 `ФВ` 号排好，把稿子里提到而硬盘里没有的图挑出来，列成一张「缺件表」。列完，缺四十一处。

那四十一处后来补齐了三十九处。有两处至今没有：一处是背面某片低地的原片，作者说他自己也没见过；一处是某张表的第七张特写，据说拍过，翻遍了没有。正文里凡涉及这两处的地方，照译，没有替他含糊过去。

体例上只有三条要交代。引《战争的雷霆》与《月之都——〈广寒宫残简〉辑佚与考释》不注页码，只注书名与章次：两书版次不同，页码会变，书名不会。原著正文用「笔者」，序言与结语用「我」；日文里这两处区别很小，作者初稿一律写「私」，我们建议他分开，他说这样反而更清楚。译文不加注，不替作者补材料，专名一律以与另两卷对得上的为准。

还有一条不在体例之内：这本书是两个人译的，而此前我们为另一部稿子闹翻过，各译各的，连校样都不通。此番是例外，因为材料横跨三卷，一个人译不下来。译文里凡有两处口径不一，都原样保留，读者若读到前后打架，那多半是我们各译各的，未必是译错。

\begin{flushright}
{\small 孙笑川}\quad{\small 李文亚}\\[2mm]
{\footnotesize\color{pencil} 二〇二四年十一月，于北京}
\end{flushright}

\clearpage

% ===================================================================
\chapter*{著者中文版序言}
\phantomsection
\addcontentsline{toc}{chapter}{著者中文版序言}
\markboth{著者中文版序言}{著者中文版序言}

我没去过月面。这一点先讲清楚。

我见过的只有照片、一卷缩微胶片，和巴兰尼科夫教授书里那几张重绘图。所以本书里凡是讲「背面长什么样」的地方，读者都可以先打个折。

校样是这么到我手上的。二〇二四年二月，事务所收到一只包裹，寄件人一栏是一个财团的名字，品名一栏写着「同人志样本」。这样的事每年有几十回，我照例拆开，里面是一叠装订成册的校样，俄文，封面是丛书页。第一页夹着一张便条，手写，只有一句「您可能想看看这个」，没有署名；翻过来，背面还有一行小字：「地址是九〇年代留下的，一直没注销。」

那天下午我把图翻了一遍。字看不懂，我把正文一页一页拍下来，丢进翻译软件，出来的日文半通不通，得对着图猜。就这个样子看了三天，看到一半才明白，我不是在读别人的东西。又过了三天，我订了机票。

二〇二四年三月我到莫斯科去看那批东西。去之前准备了三本笔记本，回来只用掉半本。东西不少，是我认不出来：照片上的地形我看得懂，照片底下那行俄文小字我看不懂。等那一叠校样译出来，已经是两个月以后。

还有一件更要紧的事。那些东西当年似乎是新闻：报纸上登过，展览办过，大人排过队。我小时候大概也看过，现在却一个字也想不起来，只剩一句「好像上去过一次」。

另一头更麻烦：这些东西后来进了我的作品。我做过几十年，把月亮画成过很多样子，画里那座城有名字、有居民、有一套自己的规矩，读者也认得。认得的方式是把它当成一部作品的一部分。我当初做的时候，只觉得这样有意思；现在回头看，那大概是我把记得的东西换了个讲法，讲得好看些，好让人愿意听下去。我做了三十年游戏，画过很多月亮，没有一张画对过。问题出在我从一开始就搞错了它在哪儿、有多大、里面有什么。

书名用的是「不想」两个字。他们在月面上确实不会变化，一落地就会了，而且很快。整本书写的是一个选择如何失效，而不是某种能力如何消失。第一次看校样时我担心这题目会被读成寓言，它没有那个意思。

书里每一个数字都有出处，对不上的都注明了。前半部分我核得动，写得细；后半部分只能推。读者若觉得前面太慢、后面太赶，原因就在这里。有一处我说得比材料能支持的要肯定，这一点在结语里认了。

\begin{flushright}
{\small 太田顺也}\\[2mm]
{\footnotesize\color{pencil} 二〇二四年十二月，于东京}
\end{flushright}

\clearpage
